\documentclass{article}
\usepackage[T1]{fontenc}
\usepackage{lmodern}
\usepackage{graphicx}
\usepackage{amsmath,amssymb}
\usepackage[a4paper,margin=2cm]{geometry}
\usepackage[absolute,overlay]{textpos}
\setlength{\TPHorizModule}{1mm}
\setlength{\TPVertModule}{1mm}
\usepackage[indonesian]{babel}
\usepackage{hyperref}
\setlength{\emergencystretch}{2em}
% unit: Halaman 20

\begin{document}

% === Halaman 20 (llm) ===

Selanjutnya untuk setiap putaran 1 sampai $r$ dilakukan operasi $XOR$, pergeseran ke kiri secara sirkuler, dan penjumlahan dalam modulo $2^{32}$ dengan kunci putaran sebagai berikut:

\begin{align*}
\text{for } i \leftarrow 1 \text{ to } r \text{ do}\
\quad A \leftarrow ((A \oplus B) \lll B) + S[2i]\
\quad B \leftarrow ((B \oplus A) \lll A) + S[2i+1]\
\text{endfor}
\end{align*}

\end{document}
